\chapter{Independence and Convergence}

\section{Independence}

Independence captures the intuitive idea that the outcome of one event does not influence the likelihood of another.

\begin{definition}[Independence of Events]
Events $(A_{\lambda})_{\lambda \in \Lambda}$ are \textbf{mutually independent} if for every finite sub-collection with indices $\lambda_1, \dots, \lambda_k$, the probability of the intersection is the product of the probabilities:
\[
\mathbb{P}\left( \bigcap_{j=1}^k A_{\lambda_j} \right) = \prod_{j=1}^k \mathbb{P}(A_{\lambda_j}).
\]
\end{definition}
\textbf{Warning}: Pairwise independence does NOT imply mutual independence. (Example: Bernstein's counterexample).

\begin{definition}[Independence of $\sigma$-algebras]
Sub-$\sigma$-algebras $\mathcal{G}_1, \dots, \mathcal{G}_n \subset \mathcal{F}$ are independent if for any choice of events $A_i \in \mathcal{G}_i$:
\[
\mathbb{P}\left( \bigcap_{i=1}^n A_i \right) = \prod_{i=1}^n \mathbb{P}(A_i).
\]
\end{definition}

\begin{definition}[Independence of Random Variables]
Random variables $(X_{\lambda})$ are independent if their generated $\sigma$-algebras $\sigma(X_{\lambda}) = X_{\lambda}^{-1}(\mathcal{G})$ are independent.
\end{definition}

\section{Modes of Convergence}
In analysis, a sequence of numbers converges to a limit. In probability, since the random outcomes depend on $\omega$, there are multiple ways to define ``closeness'' to a limit random variable $Y$.

Let $(Y_n)_{n \geq 1}$ and $Y$ be random variables.

\subsection{Almost Sure Convergence (a.s.)}
Pointwise convergence for ``almost all'' outcomes.
\[
Y_n \xrightarrow{\text{a.s.}} Y \iff \mathbb{P}\left( \{ \omega : \lim_{n \to \infty} Y_n(\omega) = Y(\omega) \} \right) = 1.
\]
This is the strongest form of pointwise convergence.

\paragraph{Decoding the Notation:}
The equation $\mathbb{P}(\{ \omega : \lim_{n\to\infty} Y_n(\omega) = Y(\omega) \}) = 1$ can be understood step-by-step:
\begin{itemize}
    \item \textbf{The Sequence $Y_n$:} The subscript $n$ does not mean a "set" of $n$ variables, but acts as an index for an infinite sequence of random variables $Y_1, Y_2, \dots$ evolving over "time" or "steps". 
    \item \textbf{Fixing a Scenario $\omega$:} If we fix a specific outcome $\omega$ (a single "universe" or deterministic scenario), each $Y_n(\omega)$ becomes just a standard real number. Thus, for a fixed $\omega$, the expression $(Y_1(\omega), Y_2(\omega), \dots)$ is a regular sequence of real numbers as seen in standard calculus.
    \item \textbf{The Limit:} The equation $\lim_{n \to \infty} Y_n(\omega) = Y(\omega)$ simply checks if this specific sequence of numbers converges to the target number $Y(\omega)$ for that chosen $\omega$.
    \item \textbf{The Probability:} We define the \textbf{set of all "good" outcomes} $G = \{ \omega \in \Omega : \lim_{n\to\infty} Y_n(\omega) = Y(\omega) \}$. The condition states that the \textbf{probability of this set $G$ is 1} (i.e., convergence happens almost surely).
    \item Equivalently, the set of "bad" outcomes where the limit diverges or goes to a different value has probability $0$. (e.g., throwing infinite consecutive "Heads" in a fair coin toss exists as a possible scenario $\omega$, but its probability is purely $0$).
\end{itemize}

\subsection{\texorpdfstring{$L^p$}{Lp} Convergence}
Convergence in the mean of the $p$-th power difference.
\[
Y_n \xrightarrow{L^p} Y \iff \|Y_n - Y\|_p = \left( \mathbb{E}[|Y_n - Y|^p] \right)^{1/p} \to 0.
\]
Implies $L^q$ convergence for any $q < p$ (by Hölder's inequality).

\paragraph{Why is this useful?}
\begin{itemize}
    \item \textbf{Moment Convergence}: It implies that the moments converge: $\mathbb{E}[Y_n] \to \mathbb{E}[Y]$ and $\text{Var}(Y_n) \to \text{Var}(Y)$. (Note: a.s. convergence does NOT guarantee this, e.g., ``mass escaping to infinity'').
    \item \textbf{Stability}: It controls the ``energy'' of the error. In finance, $L^2$ convergence is crucial because the Itô Integral is constructed as a limit in $L^2$ space.
    \item \textbf{Completeness}: The space $L^p$ is complete (Banach space). This means that if a sequence of random variables gets closer to each other (Cauchy), it is guaranteed to converge to a valid random variable within the space.
\end{itemize}

\subsection{Convergence in Probability}
The probability of observing a significant difference goes to zero.
\[
Y_n \xrightarrow{\mathbb{P}} Y \iff \forall \epsilon > 0, \quad \mathbb{P}(|Y_n - Y| > \epsilon) \to 0.
\]
\textbf{Relations}:
\begin{itemize}
    \item $L^p \implies$ Probability (Markov inequality).
    \item Almost Sure $\implies$ Probability.
    \item The converses are generally false (though Probability $\implies$ a.s. along a subsequence).
\end{itemize}

\subsection{Convergence in Law (Distribution)}
Only the distributions approach the limit distribution. The random variables themselves might not be close.
\[
Y_n \xrightarrow{\text{law}} Y \iff \mathbb{E}[f(Y_n)] \to \mathbb{E}[f(Y)] \quad \text{for all bounded continuous } f.
\]
Also denoted as $\implies$ or $\xrightarrow{\mathcal{D}}$. This is the weakest form of convergence.

\section{Limit Theorems}

Let $(X_n)$ be a sequence of independent and identically distributed (i.i.d.) random variables.

\begin{theorem}[Strong Law of Large Numbers]
If $X_1 \in L^1$ (measurable mean $m = \mathbb{E}[X_1]$), then the empirical average converges almost surely to the theoretical mean:
\[
\bar{X}_n = \frac{1}{n} \sum_{i=1}^n X_i \xrightarrow{\text{a.s.}} m.
\]
\end{theorem}

\begin{theorem}[Central Limit Theorem]
If $X_1 \in L^2$ (finite variance $\sigma^2 = \text{Var}(X_1)$), then the standardized sum converges in law to a standard Gaussian:
\[
\frac{\sqrt{n}(\bar{X}_n - m)}{\sigma} = \frac{\sum_{i=1}^n (X_i - m)}{\sigma \sqrt{n}} \xrightarrow{\text{law}} \mathcal{N}(0, 1).
\]
\end{theorem}
This explains the ubiquity of the Normal distribution in finance (returns aggregation).
